\section{Planificación del Proyecto}
El plan del proyecto se compone de seis planes: alcance, cronograma, costos,
calidad, riesgos, y recursos y comunicación. El alcance ya se describió; la
calidad tiene su propia sección (sección~\ref{sec:calidad}). Aquí se
presentan los demás.

\subsection{Triple Restricción}
Alcance, tiempo y costo dependen entre sí: si una se mueve, las otras dos lo
resienten.
\begin{itemize}
	\item \textbf{Alcance:} un acceso peatonal, turno vespertino y prototipo a
	      escala. Lo que queda fuera está escrito desde el inicio.
	\item \textbf{Tiempo:} 17 semanas fijas, del 5 de octubre de 2026 al 29 de
	      enero de 2027. El calendario escolar no se mueve.
	\item \textbf{Costo:} cerca de 4,100 pesos aportados por el propio equipo,
	      sin financiamiento externo, con una reserva de 10\,\%.
\end{itemize}
La regla del equipo es que, si algo se atrasa, se cede alcance antes que fecha
o calidad.

\subsection{Estrategias Generales}
Seis decisiones guían el proyecto:
\begin{itemize}
	\item \textbf{Empezar pequeño:} un acceso, un turno y un prototipo a escala.
	      Primero evidencia barata; escalar solo si los números la respaldan.
	\item \textbf{Trabajar en paralelo:} el hardware y el software avanzan
	      mientras se levanta el diagnóstico; la integración se deja para el final.
	\item \textbf{Usar lo barato y lo abierto:} componentes de bajo costo y
	      herramientas libres, en lugar de biometría o torniquetes industriales.
	\item \textbf{Medir antes y después:} el mismo instrumento se aplica antes
	      y durante el piloto, para comparar con números y no con impresiones.
	\item \textbf{Guardar holgura:} reserva del 10\,\% en costo y una semana
	      libre antes de vacaciones para absorber atrasos.
	\item \textbf{No retirar nada:} la revisión manual sigue funcionando todo el
	      piloto. El sistema se suma, no sustituye a nadie.
\end{itemize}

\subsection{Enfoque y Cronograma}
El enfoque es híbrido. Es predictivo en la estructura porque hay fechas que no
se pueden mover: el permiso de la dirección es requisito para entrar al
acceso, los materiales tardan en llegar, el piloto debe terminar antes de las
vacaciones y la entrega final es el 29 de enero de 2027. Es ágil en el
desarrollo porque aún no se sabe cuánto tarda el lector en leer un código ni
cómo afectan la luz, el brillo de la pantalla o la distancia. Esa
incertidumbre exige probar, medir y ajustar. El lector, el servidor y el panel
se construyen en iteraciones de una semana: cada lunes se eligen tareas de una
lista priorizada y cada viernes se muestra lo terminado y se revisa el avance
contra el cronograma. La Tabla~\ref{tab:etapas} muestra las etapas.

\begin{table}[tbp]
	\caption{Etapas, Forma de Trabajo e Hitos}
	\label{tab:etapas}
	\small\setlength{\tabcolsep}{3pt}
	\begin{tabular}{P{2.2cm}P{3.3cm}P{6.4cm}P{3cm}}
		\toprule
		Etapa & Fechas & Forma de trabajo & Hito \\
		\midrule
		Planeación & 5 al 16 de octubre & Predictiva: permisos, aviso de privacidad y lista de materiales. & Permiso de la dirección \\
		Diagnóstico & 12 al 30 de octubre & Predictiva: encuesta, entrevistas y observación del acceso durante cinco días; compra de materiales. & Informe de diagnóstico \\
		Desarrollo & 26 de octubre al 27 de noviembre & Ágil: cinco iteraciones semanales en dos frentes paralelos, hardware y lector, y servidor, base de datos y página web. & Prototipo integrado \\
		Pruebas & 23 de noviembre al 18 de diciembre & Iterativa en el salón (23 de noviembre al 4 de diciembre) y predictiva en el piloto (7 al 18 de diciembre). & Fin del piloto \\
		Cierre & 4 al 29 de enero & Predictiva: comparación de resultados, informe final y manual de uso. & Entrega final \\
		\bottomrule
	\end{tabular}
	\tablenote{Fuente: Ficha Inicial del Proyecto (Entregable 1). Las semanas del 21 y del 28 de diciembre son vacaciones de invierno y no tienen trabajo programado. Las tareas se registran en un tablero compartido y el código en un repositorio Git.}
\end{table}

Las actividades corren en cuatro frentes en paralelo y una integración al
final (Tabla~\ref{tab:frentes}). La integración se hace en el salón y, hasta
que pasa ahí, el prototipo no se instala en el acceso.

\begin{table}[tbp]
	\caption{Frentes de Trabajo}
	\label{tab:frentes}
	\small\setlength{\tabcolsep}{3pt}
	\begin{tabular}{P{5cm}P{2.6cm}P{7cm}}
		\toprule
		Frente y responsable & Semanas & Resultado \\
		\midrule
		Hardware y torniquete (Daniel) & 3 a 6 & Prototipo a escala con el lector montado. \\
		Lector y firma del código (Ian y Miguel) & 3 a 7 & Lector con cámara y código que caduca en segundos. \\
		Servidor, datos y panel (Juan Fernando) & 3 a 8 & Validación, registro y consultas del personal. \\
		Campo y pruebas (Bryan) & 2 a 5 y 10 a 11 & Diagnóstico, pruebas y piloto. \\
		Integración en el salón & 8 y 9 & Prototipo completo probado fuera del acceso. \\
		\bottomrule
	\end{tabular}
	\tablenote{Fuente: presentación del avance. La semana 1 es la del 5 de octubre de 2026.}
\end{table}

\subsection{Costos}
El presupuesto es una estimación preliminar del equipo, hecha con precios de
referencia del mercado mexicano. Se confirma al cotizar (Tabla~\ref{tab:costos}).

\begin{table}[tbp]
	\caption{Presupuesto Estimado}
	\label{tab:costos}
	\small\setlength{\tabcolsep}{3pt}
	\begin{tabular}{P{8cm}r}
		\toprule
		Categoría & Estimado (MXN) \\
		\midrule
		Hardware del prototipo & 1,750 \\
		Servidor, dominio y hosting (4 meses) & 970 \\
		Impresiones y material de campo & 600 \\
		Traslados y compras & 400 \\
		Reserva de contingencia (10\,\%) & 380 \\
		\midrule
		Total estimado & 4,100 \\
		\bottomrule
	\end{tabular}
	\tablenote{Fuente: presentación del avance. Las cerca de 800 horas de trabajo del equipo no se monetizan. Si la escuela presta el servidor, el total baja 970 pesos.}
\end{table}

El control de costos sigue cuatro reglas: una línea base aprobada antes de la
primera compra, cada compra registrada con comprobante y responsable, un umbral
de variación de $\pm$10\,\% (pasarlo obliga a replantear) y la comparación del
gasto acumulado contra lo planeado en la junta semanal.

\subsection{Recursos}
Con lo siguiente cuenta el equipo para ejecutar el proyecto:
\begin{itemize}
	\item \textbf{Humanos:} seis integrantes, unas 8 horas por semana cada uno,
	      cerca de 800 horas en total.
	\item \textbf{Materiales:} lector con cámara, placa de control, motor, luces
	      y zumbador, fuente de energía, base de MDF y cableado.
	\item \textbf{Técnicos:} servidor con base de datos, repositorio de código y
	      entorno de programación.
	\item \textbf{Institucionales:} permiso de la dirección, espacio con
	      corriente y Wi-Fi, y apoyo de vigilancia.
	\item \textbf{Financieros:} aportación del equipo, cerca de 4,100 pesos.
	\item \textbf{De información:} padrón de estudiantes de Control Escolar y
	      resultados del diagnóstico.
\end{itemize}
Lo que depende de terceros (permiso, Wi-Fi, padrón) es lo que encabeza el
registro de riesgos.

\subsection{Comunicación}
El equipo se comunica con tres mecanismos: una junta semanal de 20 minutos
para revisar el avance contra el cronograma, un reporte quincenal a la
docente con lo hecho y lo pendiente, y un tablero de tareas en el repositorio
con una columna por frente. Todo cambio que toque fecha, costo o alcance lo
autoriza el líder del proyecto y queda por escrito.

\subsection{Riesgos}
\label{sec:riesgos}
El registro tiene siete riesgos. Cada uno tiene un responsable que lo vigila
y se revisa en la junta semanal. Los dos primeros condicionan la fecha del
piloto (Tabla~\ref{tab:riesgos}).

\begin{table}[tbp]
	\caption{Registro de Riesgos}
	\label{tab:riesgos}
	\small\setlength{\tabcolsep}{3pt}
	\begin{tabular}{P{4.8cm}P{2.3cm}P{1.8cm}P{5.7cm}}
		\toprule
		Riesgo & Probabilidad & Impacto & Respuesta \\
		\midrule
		La dirección no autoriza el piloto & Media & Alto & Solicitud formal en la semana 1; plan B en el laboratorio. \\
		Wi-Fi inestable en el acceso & Alta & Alto & Modo fuera de línea con lista local y sincronización. \\
		Retraso en la entrega de componentes & Media & Medio & Compra con dos semanas de anticipación y repuestos. \\
		Mala lectura con sol directo & Media & Medio & Visera, ajuste de exposición y pruebas con varias luces. \\
		Estudiantes sin celular o sin batería & Alta & Bajo & El carril manual sigue abierto durante todo el piloto. \\
		Manejo indebido de datos personales & Baja & Alto & Aviso de privacidad, datos mínimos, cifrado y borrado. \\
		Vacaciones y carga de fin de semestre & Alta & Medio & Piloto cerrado antes de las vacaciones; enero para el informe. \\
		\bottomrule
	\end{tabular}
	\tablenote{Fuente: presentación del avance. Los nombres de los responsables de cada riesgo no aparecen en las fuentes.}
\end{table}

\subsection{Cierre del Proyecto}
El 29 de enero de 2027 se hará un cierre en seis pasos: verificar cada
entregable contra su criterio de aceptación, entregar el informe final con la
comparación de los tres indicadores, firmar un acta de aceptación por la
dirección y la docente, transferir el manual y dar una capacitación breve a
vigilancia y Control Escolar, borrar o resguardar los registros según el aviso
de privacidad, y documentar las lecciones aprendidas con una recomendación
sobre llevar el sistema a los demás accesos o detenerlo.
